
\chapter{Introduction}
\section{Background}

\pack\ is a free and open-source command-driven software for gravity anomalies~\citep{gharti2018}, magnetic anomalies~\citep{gharti2019c}, coseismic and postearthquake deformation~\citep{gharti2019a}, earthquake-induced gravity perturbations~\citep{gharti2019b}, and normal modes based on the spectral-(infinite)-element method~\citep[e.g.,][]{patera1984,canuto1988,seriani1994,faccioli1997,komatitsch1998,komatitsch1999,peter2011,gharti2017a}. The software can run on a single processor as well as multi-core machines or large clusters. It is written mainly in FORTRAN 90, and parallelized using
MPI~\citep{gropp1994,pacheco1997} based on domain decomposition. For the domain decomposition, the open-source graph partitioning library SCOTCH~\citep{pellegrini1996} is used. The element-by-element preconditioned conjugate-gradient method~\citep[e.g.,][]{hughes1983,law1986,king1987,barragy1988} is implemented to solve the linear equations. For elastoplastic failure,
a Mohr-coulomb failure criterion is used with a viscoplastic strain method~ \citep{zienkiewicz1974}.\\

The software currently does not include an inbuilt mesher. Existing tools, such as Gmsh~\citep{geuzaine2009}, CUBIT~\citep{cubit2011}, TrueGrid~\citep{truegrid2006}, etc., can be used for hexahedral meshing, and the resulting mesh file can be converted to the input files required by \pack. Output data can be visualized and processed using the open-source visualization application ParaView (\texttt{www.paraview.org}).

\section{Status summary}
\begin{desclist}{Earthquake-induced gravity perturbation}
\item[Gravity anomalies]        : Yes
\item[Magnetic anomalies]           : Yes
\item[Coseismic deformation]                 : Yes
\item[Postearthquake deformation]                 : Yes [Experimental]
\item[Earthquake-induced gravity perturbation]                    : Yes [Experimental]
\item[Frequency-domain method]: Yes [Experimental]
\end{desclist}

\section*{Revisions}

HNG, May 12, 2022; HNG, Jan 12, 2012; HNG, Sep 08, 2011; HNG, Jul 12, 2011; HNG, May 20, 2011; HNG, Jan 17, 2011


